% Pierwsza prezentacja, SirenCodec. Kompilacja: tectonic (+ BibTeX).
% Zakres slajdów 3--9 = 5 podpunktów z polecenia (kolejność jak w ćwiczeniu).
\documentclass[aspectratio=169,11pt]{beamer}

\usepackage{fontspec}
\usepackage[polish]{babel}
\usepackage[
  backend=bibtex,
  style=numeric,
  sorting=none,
  maxbibnames=99,
  doi=true,
  url=true,
  isbn=false,
  abbreviate=false,
]{biblatex}
\addbibresource{pierwsza_prezentacja_slide.bib}
\renewcommand*{\bibfont}{\fontsize{5.2pt}{6.1pt}\selectfont}
\setbeamertemplate{bibliography item}{\insertbiblabel}

\setsansfont{lmsans10}[
  Extension      = .otf,
  UprightFont    = *-regular,
  BoldFont       = *-bold,
  ItalicFont     = *-oblique,
  BoldItalicFont = *-boldoblique,
  Ligatures        = TeX,
  Scale            = MatchLowercase,
]
\setmainfont{lmroman10}[
  Extension      = .otf,
  UprightFont    = *-regular,
  BoldFont       = *-bold,
  ItalicFont     = *-italic,
  BoldItalicFont = *-bolditalic,
  Ligatures        = TeX,
]

\usetheme{default}
\usecolortheme{default}
\usefonttheme{professionalfonts}
\renewcommand{\familydefault}{\sfdefault}

\setbeamercolor{titlebar}{fg=black,bg=gray!25}

\setbeamertemplate{navigation symbols}{}
\setbeamertemplate{footline}{%
  \hfill\insertframenumber{}/\inserttotalframenumber\hspace*{0.5em}\vspace*{0.3em}%
}

\title{Ultra niskobitowa kompresja mowy}
\subtitle{\small Seminarium magisterskie · praca dyplomowa (2 autorów)}
\author{Kosma Gąsiorowski \and Jakub Aszyk}
\institute{Politechnika Poznańska\\[0.25em]
  {\small Kierunek: Informatyka \quad Przedmiot: Internet}\\[0.35em]
  Promotor: Mariusz Nowak}
\date{\today}

\newcommand{\exerciseheader}{%
  \begin{center}\normalsize Ćwiczenie Pierwsza prezentacja\end{center}\vspace{0.5em}}

\begin{document}
\nocite{*}

\begin{frame}[plain]
  \titlepage
\end{frame}

% --- Slajd 2: jak w materiale z ćwiczenia ---
\begin{frame}[plain]{}
  \exerciseheader
  \begin{beamercolorbox}[wd=\paperwidth,sep=0.55em,center]{titlebar}
    \centering\Large\bfseries Pierwsza prezentacja
  \end{beamercolorbox}
  \vspace{0.55em}
  {\small
  \begin{itemize}\setlength{\itemsep}{0.35em}
    \item Około 10 minut
    \item Cel: \textbf{Uzasadnienie wyboru tematu i realności jego wykonania w ramach pracy dyplomowej}
    \item Zakres (sl.\ 3--9):
      \begin{itemize}\setlength{\itemsep}{0.22em}
        \item Temat, promotor, zadania szczegółowe
        \item Planowane rezultaty
        \item Literatura
        \item Harmonogram prac (punkty kontrolne, terminy)
        \item Czynniki ryzyka
      \end{itemize}
  \end{itemize}
  }
\end{frame}

% --- (1) Temat, promotor, zadania szczegółowe ---
\begin{frame}{Temat, promotor, zadania szczegółowe}
  \textbf{Temat:} \inserttitle\\[0.35em]
  \textbf{Realizacja:} projekt \textbf{SirenCodec} (kodek neuronowy, niski bitrate).\\[0.5em]
  \textbf{Promotor:} Mariusz Nowak \quad \textbf{Autorzy:} Kosma Gąsiorowski, Jakub Aszyk\\[0.55em]
  \textbf{Zadania szczegółowe:}
  \begin{enumerate}
    \item Przegląd literatury i metryk.
    \item Implementacja, trening, eksperymenty (m.in.\ od ok.\ pół roku pracy badawczej).
    \item Rozdział w pracy: metoda, wyniki, wnioski.
  \end{enumerate}
\end{frame}

% --- (2) Planowane rezultaty ---
\begin{frame}{Planowane rezultaty}
  \begin{itemize}
    \item Działający łańcuch: nagranie $\rightarrow$ kody dyskretne $\rightarrow$ odtworzenie; szacowany bitrate.
    \item Zestaw wyników: logi, checkpointy, przykłady audio, wykresy do rozdziału.
    \item Odtwarzalność eksperymentów (środowisko, komendy).
  \end{itemize}
\end{frame}

% --- (3) Literatura ---
\begin{frame}[t,shrink=5]{Literatura (1/2)}
  \vspace{-0.35em}
  \printbibliography[heading=none,keyword=litA]
\end{frame}

\begin{frame}[t,shrink=5]{Literatura (2/2)}
  \vspace{-0.35em}
  \printbibliography[heading=none,keyword=litB]
\end{frame}

% --- (4a) Harmonogram --- terminy ---
\begin{frame}{Harmonogram prac --- terminy}
  \begin{itemize}
    \item XII 2025\,--\,I 2026: przegląd literatury, przygotowanie środowiska.
    \item II 2026: implementacja (kod, dane).
    \item III\,--\,VI 2026: eksperymenty i pisanie pracy równolegle.
    \item VII 2026: obrona.
  \end{itemize}
\end{frame}

% --- (4b) Harmonogram --- punkty kontrolne ---
\begin{frame}{Harmonogram prac --- punkty kontrolne}
  \begin{itemize}
    \item Zamknięcie przeglądu literatury i gotowość środowiska (koniec I\,2026).
    \item Pierwszy stabilny trening + próbki audio (II\,2026).
    \item Domknięcie pomiarów; konsolidacja rozdziału wyników (VI\,2026).
    \item Wersja pracy do akceptacji promotora przed obroną (VII\,2026).
  \end{itemize}
\end{frame}

% --- (5) Czynniki ryzyka ---
\begin{frame}{Czynniki ryzyka}
  \begin{itemize}
    \item Jakość przy bardzo niskim bitrate (artefakty) --- ewentualna korekta celów po uzgodnieniu z promotorem.
    \item Czas obliczeń na GPU --- krótsze przebiegi, podzbiór danych.
    \item Koordynacja dwóch autorów --- podział modułów i gałęzi w repozytorium.
  \end{itemize}
\end{frame}

\begin{frame}[plain]
  \begin{center}
    {\Large Dziękujemy za uwagę.}\\[1.2em]
    Pytania?
  \end{center}
\end{frame}

\end{document}
